\par\medskip\noindent\textbf{An independent calculation of the generic torus torsor.}
Here is a direct verification of the required vanishing
for the actual cyclic rank-two torus. Put
$K=\mathbb C(B/Q)$. The splitting action on its
character lattice has order $d\in\{1,2,3,4,6\}$;
if the base group is larger, pass first to the fixed
field of the kernel of this lattice action. This gives
the actual cyclic splitting extension $L/K$ of degree $d$.
We prove
\[
                         H^1(K,T)=0.                 \tag{C.EO.H1}
\]
The cohomology here is the Galois cohomology of the
algebraic torus obtained by descent of the two-dimensional
split multiplicative group, and classifies its actual
torsors over $K$.

First every function field $K'$ of a smooth compact
complex curve has the following elementary property:
a homogeneous polynomial of degree $a$ in $b>a$
variables has a nonzero zero over $K'$. To verify it,
let $g$ be the genus of the curve, choose a point $P$,
and choose an effective divisor $D$ bounding the poles
of all coefficients. Seek each variable in
$H^0(\OO(rP))$. For large $r$ there are
$b(r+1-g)$ scalar unknowns, whereas the polynomial
value belongs to $H^0(\OO(arP+D))$, of dimension
$ar+\deg D+1-g$. Its vanishing imposes at most that
many homogeneous equations of positive degree.
For large $r$ the number of unknowns is strictly
larger. The height bound for an ideal generated by
that many elements gives an affine homogeneous zero
set of dimension at least one. Over $\mathbb C$ it
contains a point other than the origin. This yields
the asserted nonzero zero over $K'$.

Consequently, for every finite field extension
$L'/K'$ the norm map is surjective:
\[
               N_{L'/K'}:(L')^*\longrightarrow(K')^*.
                                                               \tag{C.EO.H2}
\]
Indeed, for $c\in(K')^*$, choose a $K'$-basis of
$L'$ of size $a=[L':K']$ and apply the preceding
property to
$N_{L'/K'}(x_1e_1+\cdots+x_ae_a)-cz^a$.
This has degree $a$ in $a+1$ variables.
A nonzero zero cannot have $z=0$, because the norm
of a field element vanishes only for the zero element.
Dividing by $z$ gives an element of norm $c$.
Every intermediate field used below is again the
function field of a smooth compact complex curve,
so the same proof applies to it.

The elliptic lattice action has, after an integral
choice of basis, the following matrices:
\[
\begin{array}{c|c}
d&M_d\\ \hline
1&I_2\\
2&-I_2\\
3&\begin{pmatrix}0&-1\\1&-1\end{pmatrix}\\
4&\begin{pmatrix}0&-1\\1&0\end{pmatrix}\\
6&\begin{pmatrix}0&-1\\1&1\end{pmatrix}.
\end{array}                                                   \tag{C.EO.H3}
\]
For completeness, a rank-two lattice preserved by a
primitive automorphism of order $3,4$, or $6$ is a
rank-one module over the corresponding ring
$\mathbb Z[\zeta_d]$. After tensoring with $\mathbb Q$
and choosing a nonzero lattice vector, it becomes
a fractional ideal of that ring. The Gaussian and
Eisenstein rings are Euclidean for their usual norm,
so this ideal is principal. The same argument applies
to the dual character lattice. It gives the matrices
in (C.EO.H3); replacing a generator by its inverse
gives an integrally conjugate matrix. The order-two
case is scalar multiplication by $-1$.

If $d=1$, the torus is split and Hilbert's Theorem 90
gives (C.EO.H1). For $d=2$, it is the product of two
quadratic norm-one tori, since each character lattice
is the sign representation. For $d=3$, it is the
cubic norm-one torus. To specify this identification,
the exact sequence
\[
 1\longrightarrow R^1_{L/K}\mathbb G_m
 \longrightarrow R_{L/K}\mathbb G_m
 \xrightarrow{\ N\ }\mathbb G_m\longrightarrow1
                                                               \tag{C.EO.H4}
\]
has character quotient
$\mathbb Z[C_3]/\mathbb Z(1+g+g^2)$.
In the basis $[1],[g]$, multiplication by $g$ is
exactly $M_3$. In degree two the same quotient is
the sign lattice. Galois descent of characters
therefore gives the stated isomorphisms of tori.
The long exact sequence of (C.EO.H4) and Hilbert 90
identify its $H^1$ with $K^*/N_{L/K}L^*$, which is
zero by (C.EO.H2).

For $d=4$, let $K_2$ be the degree-two intermediate
field. The torus is
\[
        T\simeq R_{K_2/K}
                  \bigl(R^1_{L/K_2}\mathbb G_m\bigr).
                                                               \tag{C.EO.H5}
\]
This is an exact character-lattice identification:
induce the sign module of $\langle g^2\rangle$
to $\langle g\rangle$; with basis $e,ge$, the
relations $g(e)=ge$ and $g(ge)=-e$ give $M_4$.
Weil restriction and the corresponding Shapiro
map identify its $H^1$ over $K$ with the $H^1$
of the norm-one torus over $K_2$. That group is
$K_2^*/N_{L/K_2}L^*$ and vanishes by (C.EO.H2).

For $d=6$, let $K_2,K_3$ be the degree-two and
degree-three intermediate fields. On restriction
to $K_2$, a generator acts by $M_6^2$; on restriction
to $K_3$ it acts by $M_6^3=-I_2$. The integral matrix
\[
 P=\begin{pmatrix}1&-1\\0&1\end{pmatrix},\qquad
                  P^{-1}M_6^2P=M_3                  \tag{C.EO.H6}
\]
identifies the first restricted torus with the
cubic norm-one torus, and the second with a product
of quadratic norm-one tori. Thus both restricted
$H^1$ groups vanish.

We spell out the transfer used to return to $K$.
For a finite extension $K'/K$, put
$\mathcal G=\operatorname{Gal}(K^{\rm sep}/K)$,
$\mathcal H=\operatorname{Gal}(K^{\rm sep}/K')$,
and $A=T(K^{\rm sep})$, written additively for this
calculation. Restriction is restriction of continuous
one-cocycles to $\mathcal H$. Choose left coset
representatives $R$ for $\mathcal G/\mathcal H$;
write $\sigma r=r_\sigma h_{\sigma,r}$.
For a one-cocycle $c$ on $\mathcal H$ the transfer
is represented by
\[
       (\operatorname{cor}c)(\sigma)
            =\sum_{r\in R}r_\sigma\,
                         c(h_{\sigma,r}).             \tag{C.EO.H7}
\]
The cocycle identity verifies that this is a cocycle.
For $c=\operatorname{res}b$, the identity
$b(\sigma r)=b(\sigma)+\sigma b(r)
=b(r_\sigma)+r_\sigma b(h_{\sigma,r})$
shows that (C.EO.H7) differs from
$[K':K]b(\sigma)$ by the coboundary of
$\sum_{r\in R}b(r)$. Thus
$\operatorname{cor}\circ\operatorname{res}$ is
multiplication by $[K':K]$ on $H^1(K,T)$.
The restrictions to $K_2$ and $K_3$ just proved
to vanish therefore kill every class by both
two and three. Their coprimality gives (C.EO.H1).

This proves the required generic torsor vanishing
in every possible elliptic character, with the
actual restriction and norm maps specified.
The integral identities (C.EO.H3), (C.EO.H5),
and (C.EO.H6), together with the compatible
elliptic quotient and additive-kernel characters,
are independently executed in
\path{code/A_equivariant_torus.py} and its saved output.
